\documentclass[12pt,a4paper]{report}

\usepackage{fontspec}
\usepackage{polyglossia}
\setdefaultlanguage{romanian}
\usepackage[margin=2.25cm]{geometry}
\usepackage{setspace}
\usepackage{amsmath}
\usepackage{amssymb}
\usepackage{enumitem}
\usepackage{longtable}
\usepackage{array}
\usepackage{booktabs}
\usepackage{xcolor}
\usepackage{listings}
\usepackage[
    colorlinks,
    linkcolor={black},
    citecolor={black},
    urlcolor={blue}
]{hyperref}

\onehalfspacing
\setlength{\emergencystretch}{3em}
\setlength{\parskip}{0.15em}
\setlist{nosep,leftmargin=*}

\definecolor{codebackground}{rgb}{0.96,0.96,0.96}
\definecolor{codecomment}{rgb}{0.23,0.45,0.30}
\definecolor{codekeyword}{rgb}{0.10,0.20,0.55}
\definecolor{codestring}{rgb}{0.55,0.15,0.15}

\lstdefinestyle{codestyle}{
    backgroundcolor=\color{codebackground},
    basicstyle=\ttfamily\scriptsize,
    breaklines=true,
    frame=single,
    columns=fullflexible,
    keepspaces=true,
    showstringspaces=false,
    numbers=left,
    numberstyle=\tiny\color{gray},
    keywordstyle=\color{codekeyword}\bfseries,
    commentstyle=\color{codecomment},
    stringstyle=\color{codestring},
    tabsize=4
}
\lstset{style=codestyle}

\newcommand{\codechunk}[2]{%
\subsection*{Liniile #1--#2}
\lstinputlisting[language=Python,firstline=#1,lastline=#2,firstnumber=#1]{../src/models/inference.py}
}

\title{Explicarea fisierului \texttt{inference.py}}
\author{}
\date{}

\begin{document}

\maketitle
\tableofcontents
\clearpage

\chapter{Rolul fisierului \texttt{inference.py}}

Fisierul \path{src/models/inference.py} contine utilitarele prin care modelul
MaskNet este incarcat dintr-un checkpoint si aplicat pe un semnal audio. Acest
fisier este important deoarece face legatura dintre modelul antrenat si folosirea
lui in demo-uri, evaluare si sistemul real-time.

Din punct de vedere functional, fisierul are doua responsabilitati principale:

\begin{itemize}
    \item reconstruirea arhitecturii MaskNet din configuratia salvata in checkpoint;
    \item transformarea unui waveform zgomotos in waveform imbunatatit, prin STFT,
    predictia mastii si ISTFT.
\end{itemize}

Exista doua moduri de inferenta:

\begin{itemize}
    \item \textbf{masca de magnitudine}: modelul primeste \(|Y|\), prezice o masca
    \(\hat{M}\), iar reconstructia foloseste faza semnalului zgomotos;
    \item \textbf{masca complexa}: modelul primeste canale real/imaginar, prezice
    o masca complexa, iar reconstructia se face direct din STFT complex.
\end{itemize}

\chapter{Importuri}

\codechunk{1}{18}

\begin{description}
    \item[Linia 1.] Docstring-ul defineste fisierul ca un modul de utilitare
    comune pentru inferenta MaskNet.

    \item[Linia 3.] Activeaza adnotarile amanate. Acest lucru permite folosirea
    unor hint-uri de tip mai moderne, fara a forta evaluarea lor imediata la
    import.

    \item[Linia 5.] Importa \texttt{Path}, folosit pentru a accepta checkpoint-uri
    atat ca string, cat si ca obiect de cale.

    \item[Liniile 7--8.] Importa NumPy si PyTorch. NumPy este folosit pentru
    waveform-ul de intrare/iesire, iar PyTorch pentru tensorii de inferenta si
    incarcarea checkpoint-ului.

    \item[Linia 10.] Importa configuratia globala a proiectului, de unde sunt
    preluate valorile implicite pentru \texttt{N\_FFT}, \texttt{HOP\_LEN} si
    \texttt{FRAME\_LEN}.

    \item[Liniile 11--17.] Importa functiile STFT/ISTFT si operatiile pentru masti
    complexe. Acestea sunt folosite pentru transformarea dintre domeniul timp si
    domeniul timp-frecventa.

    \item[Linia 18.] Importa clasa \texttt{MaskNet}, necesara pentru reconstruirea
    modelului inainte de incarcarea ponderilor.
\end{description}

\chapter{Incarcarea checkpoint-ului}

\codechunk{21}{24}

\begin{description}
    \item[Linia 21.] Defineste functia \texttt{load\_masknet\_checkpoint}. Ea primeste
    calea catre checkpoint si dispozitivul PyTorch pe care va fi incarcat modelul.
    Tipul returnat este \texttt{MaskNet}.

    \item[Linia 22.] Docstring-ul explica scopul direct: incarcarea modelului din
    checkpoint-ul de training.

    \item[Linia 23.] \texttt{torch.load} citeste fisierul checkpoint. Parametrul
    \texttt{map\_location=device} asigura incarcarea direct pe CPU sau GPU, in
    functie de context. \texttt{weights\_only=False} permite citirea intregului
    dictionar salvat, nu doar a ponderilor.

    \item[Linia 24.] Extrage \texttt{model\_config} din checkpoint. Daca acesta nu
    exista, se foloseste dictionar gol, iar modelul va fi construit cu valori
    implicite.
\end{description}

\codechunk{26}{38}

\begin{description}
    \item[Liniile 26--38.] Reconstruiesc arhitectura MaskNet pe baza configuratiei
    salvate. Aceasta etapa este esentiala: ponderile unui model pot fi incarcate
    corect doar daca arhitectura are aceeasi structura ca la antrenare.

    \item[Canale de intrare/iesire.] \texttt{in\_channels} si \texttt{output\_channels}
    disting intre modelul de magnitudine si modelul complex. Pentru magnitudine,
    valorile sunt de obicei 1. Pentru CRM, valorile sunt 2.

    \item[Componente optionale.] \texttt{use\_lstm}, \texttt{use\_attention},
    \texttt{use\_residual} si \texttt{use\_depthwise} refac exact varianta arhitecturala
    folosita in checkpoint.

    \item[Tipul mastii.] \texttt{mask\_type}, \texttt{output\_activation},
    \texttt{output\_scale} si \texttt{complex\_mask\_clip} controleaza modul in care
    iesirea modelului este interpretata.
\end{description}

\codechunk{40}{51}

\begin{description}
    \item[Liniile 40--43.] Codul accepta doua denumiri posibile pentru ponderi:
    \texttt{model\_state\_dict} si \texttt{model\_state}. Aceasta compatibilitate
    este utila daca diferite scripturi de training au salvat checkpoint-urile cu
    chei usor diferite.

    \item[Liniile 44--47.] Daca checkpoint-ul nu contine niciuna dintre cheile
    asteptate, se ridica o eroare explicita. Mesajul include cheile gasite, ceea
    ce ajuta la diagnosticarea unui checkpoint incompatibil.

    \item[Linia 49.] Muta modelul pe dispozitivul cerut: CPU sau CUDA.

    \item[Linia 50.] Pune modelul in modul \texttt{eval}. Acest lucru dezactiveaza
    comportamentele de training, precum dropout-ul, si face BatchNorm-ul sa foloseasca
    statisticile invatate.

    \item[Linia 51.] Intoarce modelul pregatit pentru inferenta.
\end{description}

\chapter{Functia principala de enhancement}

\codechunk{54}{69}

\begin{description}
    \item[Liniile 54--61.] Defineste functia \texttt{enhance\_audio\_with\_masknet}.
    Ea primeste waveform-ul zgomotos ca \texttt{np.ndarray}, modelul incarcat,
    dispozitivul PyTorch si, optional, parametrii STFT.

    \item[Linia 62.] Docstring-ul formuleaza scopul: imbunatatirea unui waveform
    zgomotos folosind modelul MaskNet antrenat.

    \item[Linia 63.] Converteste semnalul NumPy in tensor PyTorch \texttt{float}
    si il muta pe dispozitivul de inferenta.

    \item[Liniile 64--66.] Daca parametrii STFT nu sunt furnizati explicit, sunt
    folositi cei din \texttt{config.py}: \texttt{N\_FFT}, \texttt{HOP\_LEN} si
    \texttt{FRAME\_LEN}.

    \item[Linia 68.] Citeste tipul mastii din model. Daca modelul nu are atributul,
    se presupune masca de magnitudine.

    \item[Linia 69.] Activeaza autocast doar pe CUDA. Aceasta permite mixed precision
    in inferenta GPU, reducand costul computational.
\end{description}

\chapter{Ramura pentru masca complexa}

\codechunk{71}{91}

\begin{description}
    \item[Linia 71.] Verifica daca modelul este de tip \texttt{complex}. In acest
    caz, modelul va lucra cu spectrograma complexa, nu doar cu magnitudinea.

    \item[Liniile 72--78.] Calculeaza STFT complex:
    \[
    Y[k,m] \in \mathbb{C}.
    \]
    Parametrul \texttt{return\_complex=True} face ca functia sa intoarca direct
    tensorul complex.

    \item[Liniile 79--82.] Inferenta ruleaza fara gradienti, prin
    \texttt{torch.inference\_mode}. In interior, autocast-ul este activ doar pe
    CUDA. \texttt{complex\_to\_channels} transforma STFT-ul complex in doua canale:
    partea reala si partea imaginara. \texttt{unsqueeze(0)} adauga dimensiunea de
    batch.

    \item[Linia 84.] Aplica masca complexa prezisa peste STFT-ul zgomotos:
    \[
    \hat{S}[k,m] = \hat{M}_{\mathbb{C}}[k,m] \cdot Y[k,m].
    \]

    \item[Liniile 85--91.] Reconstruieste waveform-ul prin ISTFT complexa. Parametrul
    \texttt{length=len(noisy\_audio)} forteaza iesirea sa aiba aceeasi lungime ca
    semnalul de intrare.
\end{description}

\chapter{Ramura pentru masca de magnitudine}

\codechunk{92}{112}

\begin{description}
    \item[Linia 92.] Ramura \texttt{else} trateaza cazul standard: masca de magnitudine.

    \item[Liniile 93--98.] Calculeaza STFT-ul si separa magnitudinea si faza:
    \[
    Y[k,m] = |Y[k,m]| e^{j\angle Y[k,m]}.
    \]

    \item[Liniile 99--102.] Ruleaza modelul fara gradienti si cu autocast optional.
    \texttt{unsqueeze(0).unsqueeze(0)} adauga dimensiunile de batch si canal,
    transformand spectrograma \((F,T)\) in \((1,1,F,T)\).

    \item[Linia 104.] Masca prezisa este eliminata de dimensiunile artificiale si
    inmultita cu magnitudinea zgomotoasa:
    \[
    |\hat{S}[k,m]| = \hat{M}[k,m] \cdot |Y[k,m]|.
    \]

    \item[Liniile 105--112.] Reconstruieste waveform-ul prin ISTFT, folosind
    magnitudinea imbunatatita si faza semnalului zgomotos. Aceasta este abordarea
    clasica pentru modelele care estimeaza doar magnitudinea.
\end{description}

\codechunk{113}{113}

\begin{description}
    \item[Linia 113.] Rezultatul final este detasat din graful PyTorch, mutat pe
    CPU si convertit inapoi in \texttt{np.ndarray}. Astfel, functia poate fi folosita
    usor in codul DSP, evaluare, demo offline si real-time.
\end{description}

\chapter{Concluzie}

\texttt{inference.py} este fisierul care transforma modelul antrenat intr-un
instrument utilizabil. \texttt{load\_masknet\_checkpoint} reconstruieste arhitectura
si incarca ponderile, iar \texttt{enhance\_audio\_with\_masknet} executa pipeline-ul
de inferenta: waveform \(\rightarrow\) STFT \(\rightarrow\) predictie masca
\(\rightarrow\) aplicare masca \(\rightarrow\) ISTFT.

Importanta fisierului este data de faptul ca aceeasi logica este reutilizata in
mai multe locuri ale proiectului: evaluare, demo offline si demo real-time.
Astfel, comportamentul modelului ramane consecvent indiferent daca este folosit
pe fisiere sau pe blocuri audio live.

\end{document}
